

\usepackage{microtype}
\usepackage{amsfonts}

\usepackage{tikz-cd}
\usepackage{amssymb}
\usepackage{enumitem}
\usepackage{mathrsfs}
\usepackage{tikz}
\usetikzlibrary{topaths}
\usetikzlibrary{calc}
\renewcommand{\theenumi}{\roman{enumi}}
\renewcommand{\labelenumi}{\rm (\theenumi)}
\usepackage{empheq}

\newcommand{\listlabel}[1]{
  \item[\quad(#1)]\hypertarget{#1}{}\global\expandafter\def\csname #1\endcsname{\hyperlink{#1}{(#1)}}}

\newcommand{\N}{\mathbb{N}}
\newcommand{\Z}{\mathbb{Z}}
\newcommand{\Q}{\mathbb{Q}}
\newcommand{\R}{\mathbb{R}}

\newcommand{\Ab}{\mathrm{Ab}}
\newcommand{\Aut}{\mathrm{Aut}}
\newcommand{\Homeo}{\mathrm{Homeo}}
\newcommand{\PSL}{\mathrm{PSL}}

\newcommand{\F}{\mathrm{F}}
\newcommand{\Supp}{\mathrm{Supp}}
\newcommand{\Sub}{\mathrm{Sub}}
\newcommand{\Stab}{\mathrm{Stab}}
\newcommand{\Fix}{\mathrm{Fix}}
\newcommand{\id}{\mathrm{id}}

\newcommand{\Hyp}{\mathcal{H}}
\newcommand{\Cay}{\mathrm{Cay}}
\newcommand{\Conf}{\mathrm{Conf}}

\newcommand{\LL}{\mathrm{LL}}
\newcommand{\Mov}{\mathrm{Mov}}
\newcommand{\Fnj}{\mathrm{F}_n^{\psi_j}}
\newcommand{\X}{\mathcal{X}}












































\theoremstyle{plain}
\newtheorem{cit}[cdrthm]{Citation}

\newtheorem{thmA}{Theorem}
\renewcommand{\thethmA}{\Alph{thmA}}

\section*{Introduction}

A recent, fruitful line of inquiry in geometric group theory involves classifying all possible cobounded hyperbolic actions of a given group. An action of a group by isometries on a metric space is called \emph{hyperbolic} if the metric space is hyperbolic, in the sense of Gromov. In order to avoid pathologies, it is standard to only consider actions that are \emph{cobounded}, meaning that the space can be covered by $G$-translates of some fixed bounded subset. There is a natural notion of order on such actions, which induces an equivalence relation: the resulting quotient set inherits the structure of a poset, called the \emph{poset of hyperbolic structures}, denoted by $\Hyp(G)$.

This poset was introduced in~\cite{ABO}, where its properties for general groups were studied, and it was shown to be typically unwieldy for groups exhibiting negative curvature, specifically acylindrically hyperbolic groups. On the other hand, for certain solvable groups, the poset of hyperbolic structures is more reasonable, and can be computed explicitly~\cite{caprace15, sahana:wreath, AR:BS, confining, ABR2}. In another direction, the poset is \emph{trivial} for higher-rank lattices~\cite{haettel} and many groups of dynamical origin acting on the circle or the Cantor set~\cite{anthony, NL}.

Our focus in this paper is another group of dynamical origin, which instead acts on the interval, and shares some of the algebraic properties of the solvable groups whose poset of hyperbolic structure has been previously studied: \emph{Thompson's group $F$}.
This is the group of all piecewise linear orientation preserving homeomorphisms of the interval $[0,1]$ with all slopes powers of 2 and breakpoints in $\Z[1/2]$. This seemingly ad hoc definition belies the rich structure and bizarre properties of $F$; for instance it was the first known torsion-free group to admit a classifying space with every skeleton finite, but not admit any finite classifying space~\cite{brown84}. It is also an example of a concrete, straightforward-to-define group for which properties like amenability, $C^*$-simplicity, and automaticity remain open. Background material on $F$ can be found for example in~\cite{belk04,cannon96}.

There is an obvious generalization of $F$ to an infinite family of groups $F_n$ ($n\ge 2$), with $F=F_2$, namely, $F_n$ is the group of all piecewise linear orientation preserving homeomorphisms of the interval $[0,1]$ with all slopes powers of $n$ and breakpoints in $\Z[1/n]$. These groups were first officially investigated by Brown in~\cite{brown87}, and are the natural ``$F$-like'' version of the Higman--Thompson groups $V_n$ introduced by Higman in~\cite{higman74}.

Our goal in this paper is to understand the poset of hyperbolic structures on $F_n$. We analyze the structure of this poset from both a ``global'' and a ``local'' perspective, first breaking the poset into large subposets, revealing a global structure, and then inspecting these smaller pieces individually to obtain the local structure. Our first main result shows that the global structure is as simple as one might expect; see Subsection~\ref{ssec:general_hyp_actions} for the relevant definitions.

\begin{thmA}[Global structure of $\Hyp(F_n)$]
\label{intro:thm:global}
The poset $\Hyp(F_n)$ has the following structure.
\begin{enumerate}
   \item The poset of general type structures $\Hyp_{gt}(F_n)$ is empty.
   \item The poset of quasi-parabolic structures $\Hyp_{qp}(F_n)$ consists of two incomparable isomorphic posets $\X_0$ and $\X_1$.
   \item The poset of lineal structures $\Hyp_{\ell}(F_n)$ coincides with the poset of oriented lineal structures $\Hyp_{\ell}^+(F_n)$, which is an uncountable antichain, with two distinguished elements $[\pm\chi_0]$ and $[\pm\chi_1]$ dominated by $\X_0$ and $\X_1$, respectively.
   \item The poset of elliptic structures $\Hyp_e(F_n)$ is reduced to the smallest element, which is dominated by every lineal structure.
\end{enumerate}
\end{thmA}

See Figure~\ref{intro:fig:global} for a picture.

\begin{figure}[htb]
\centering
\begin{tikzpicture}[line width=0.8pt]

\draw (0,2) -- (0,0) -- (4,-2) -- (8,0) -- (8,2);
\draw (2,0) -- (4,-2) -- (4,0)   (4,-2) -- (6,0);

\filldraw[white] (4,-2) circle (8pt);
\draw (4,-2) circle (8pt);
\node at (4,-2) {$*$};
\node at (-2.5,-2) {$\Hyp_e(F_n)$:};

\filldraw[white] (0,0) circle (16pt);
\draw (0,0) circle (16pt);
\node at (0,0) {$[\pm\chi_0]$};

\filldraw[white] (8,0) circle (16pt);
\draw (8,0) circle (16pt);
\node at (8,0) {$[\pm\chi_1]$};
\node at (-2.5,0) {$\Hyp_\ell(F_n)$:};

\filldraw[white] (0,2) circle (18pt);
\draw (0,2) circle (18pt);
\node at (0,2) {$\mathcal{X}_0$};

\filldraw[white] (8,2) circle (18pt);
\draw (8,2) circle (18pt);
\node at (8,2) {$\mathcal{X}_1$};
\node at (-2.5,2) {$\Hyp_{qp}(F_n)$:};

\filldraw[white] (1,1) -- (7,1) -- (7,-0.3) -- (1,-0.3);
\node at (2,0) {$\cdots$};
\node at (4,0) {$\cdots$};
\node at (6,0) {$\cdots$};

\end{tikzpicture}
\caption{The global structure of the poset $\Hyp(F_n)$, as described in Theorem~\ref{intro:thm:global}.}
\label{intro:fig:global}
\end{figure}
















































































































































\section{Hyperbolic actions}\label{sec:actions}

In this section we recall some background on hyperbolic actions, the poset of hyperbolic structures, and Busemann pseudocharacters.


\subsection{General hyperbolic actions}\label{ssec:general_hyp_actions}

Let us recall the general classification of hyperbolic actions, where by a \emph{hyperbolic action} we mean an action of a group $G$ by isometries on a hyperbolic metric space $X$. An element $g\in G$ is called \begin{enumerate}
    \item \emph{elliptic} if $\langle g\rangle$ has bounded orbits, \item \emph{loxodromic} if the map $n\mapsto g^n x$ is a quasi-isometric embedding from $\Z$ to $X$ for some (equivalently every) $x\in X$, and 
    \item \emph{parabolic} if it is neither elliptic nor loxodromic. 
    \end{enumerate}
    A loxodromic element $g$ fixes precisely two points in the (Gromov) boundary $\partial X$ of $X$, namely the limits at $+\infty$ and $-\infty$ of the sequence $g^n x$ for some (any) $x\in X$. A parabolic element has exactly one fixed point in $\partial X$. 

Let $\Lambda(G)$ denote the set of limit points of $G$ in $\partial X$. More precisely, $\Lambda(G)$ is the intersection of $\partial X$ with the closure in $X\cup\partial X$ of some (equivalently any) $G$-orbit $G.x$. Now hyperbolic actions are classified as follows; this is essentially due to Gromov~\cite[Section~8.2]{gromov87}, with the level of generality due to Hamann~\cite{hamann17}, and we will use the phrasing from~\cite[Theorem~2.7]{confining}.

\begin{cit}
Let $G$ be a group acting by isometries on a hyperbolic space $X$. Then precisely one of the following holds:
\begin{enumerate}
    \item $|\Lambda(G)|=0$, in which case the action is called \emph{elliptic}.
    \item $|\Lambda(G)|=1$, in which case the action is called \emph{parabolic} (or \emph{horocyclic}).
    \item $|\Lambda(G)|=2$, in which case the action is called \emph{lineal}.
    \item $|\Lambda(G)|=\infty$, in which case the action is called \emph{non-elementary}.
\end{enumerate}
\end{cit}

Let us collect some facts about such actions. An action being elliptic is equivalent to having bounded orbits. A parabolic action cannot be cobounded. An action being lineal is equivalent to saying that at least one loxodromic element exists and all loxodromic elements share the same pair of limit points in $\partial X$. We call a lineal action \emph{oriented} if the two limit points are fixed, and \emph{non-oriented} if some element of the group can swap them. Non-elementary actions further split into the following two cases:

\begin{defi}[Quasi-parabolic, general type]
Let $G$ be a group acting non-elementarily on a hyperbolic space $X$. If $G$ fixes a (necessarily unique) point of $\partial X$, or equivalently if all loxodromic elements have one fixed point in $\partial X$ in common, then we call the action \emph{quasi-parabolic} (or \emph{focal}). If $G$ does not fix any points in $\partial X$ then we call the action \emph{of general type}.
\end{defi}

Note that if a group admits a general type action on a hyperbolic space, then a standard ping-pong argument shows that the group contains non-abelian free subgroups.


\subsection{The poset of cobounded hyperbolic actions}\label{ssec:poset}

Given a group, one can inspect all of its cobounded hyperbolic actions, and ask whether there is some structure relating them. In~\cite{ABO}, the first author together with Abbott and Osin introduced the poset $\Hyp(G)$ of (equivalence classes of) hyperbolic actions of $G$. Let us recall the construction of $\Hyp(G)$ now.

We will denote an action of $G$ on a set $X$ using the notation $G \curvearrowright X$. Given two cobounded hyperbolic actions $G\curvearrowright X$ and $G\curvearrowright Y$, we say that $G\curvearrowright X$ is \emph{dominated by} $G\curvearrowright Y$, denoted $G\curvearrowright X \preceq G\curvearrowright Y$, if there exists a coarsely Lipschitz map $f\colon Y\to X$ that is coarsely $G$-equivariant with respect to the two actions. Here $f$ is \emph{coarsely Lipschitz} if it distorts distances at most linearly, and \emph{coarsely $G$-equivariant} if for all $x\in X$ there is a uniform bound on the distance from $f(gx)$ to $gf(x)$ for all $g\in G$.

The relation $\preceq$ is clearly reflexive and transitive, so it is a preorder. We thus get an induced equivalence relation on cobounded hyperbolic actions of $G$, with two actions being equivalent if they dominate each other. We denote by $[G\curvearrowright X]$ the equivalence class of $G\curvearrowright X$, and write $\preccurlyeq$ for the partial order induced by $\preceq$ on the set of equivalence classes.

\begin{defi}[The poset $\Hyp(G)$]
We denote by $\Hyp(G)$ the poset of all equivalence classes of cobounded hyperbolic actions of $G$, with the partial order $\preccurlyeq$.
\end{defi}

It turns out that, up to equivalence, every element of $\Hyp(G)$ is represented by an action of $G$ on a (hyperbolic) Cayley graph of $G$, usually with respect to an infinite generating set. Let us set things up in the language of Cayley graphs, and then tie things together. Given any two generating sets $A$ and $B$ for $G$, say that $A$ is \emph{dominated by} $B$, denoted $A\preceq B$, if there is a uniform bound on the word length in $A$ of elements of $B$. Declare that $A$ and $B$ are \emph{equivalent} if they dominate each other. For example, if $G$ is finitely generated then all its finite generating sets are equivalent to each other. Also note that $G$ itself is dominated by all generating sets. Consider equivalence classes~$[A]$ of generating sets $A$ of $G$ using this equivalence relation, such that the Cayley graph $\Cay(G,A)$ is hyperbolic. Call such an $[A]$ a \emph{hyperbolic structure on $G$}. The hyperbolic structures on $G$ form a poset, with $[A]\preccurlyeq [B]$ whenever $A\preceq B$. For example, if $G$ is a hyperbolic group then the equivalence class of finite generating sets is the largest element of this poset, and for any $G$ the equivalence class $[G]$ is the smallest element. (We should point out that if $A$ and $B$ are generating sets with $A\subseteq B$ then $[B]\preccurlyeq [A]$, so the order is reversed in some sense.)

The following is immediate from Proposition~3.12 of~\cite{ABO} and the discussion after this proposition.

\begin{cit}
\label{cit:cayley:representative}

The poset of hyperbolic structures on $G$ is isomorphic to $\Hyp(G)$, via the poset isomorphism $[A]\mapsto [G\curvearrowright \Cay(G,A)]$.
\end{cit}

We may now view $\Hyp(G)$ either as the poset of equivalence classes of cobounded hyperbolic actions of $G$, or as the poset of hyperbolic structures on $G$, with the two viewpoints being useful for different purposes. We may also refer to a hyperbolic structure as being \emph{represented by} a certain action, with this isomorphism being implicitly understood. 

For a generating set $A$ of $G$ with $\Cay(G,A)$ hyperbolic, since the action of $G$ on $\Cay(G,A)$ is a cobounded hyperbolic action, we know it is either elliptic, lineal, quasi-parabolic, or general type. Also, if $[A]=[B]$ then the actions of $G$ on $\Cay(G,A)$ and $\Cay(G,B)$ have the same type, so $\Hyp(G)$ decomposes as a disjoint union
\[
\Hyp(G) = \Hyp_e(G) \sqcup \Hyp_\ell(G) \sqcup \Hyp_{qp}(G) \sqcup \Hyp_{gt}(G)\text{,}
\]
where these subsets denote the sets of (equivalence classes of) actions that are elliptic, lineal, quasi-parabolic, and general type, respectively. Let us also write $\Hyp_\ell^+(G)$ for the subposet of elements of $\Hyp_\ell(G)$ represented by oriented lineal actions. Clearly $[A]\in\Hyp_e(G)$ if and only if $\Cay(G,A)$ is bounded, and all such $[A]$ equal $[G]$, so this case is uninteresting: we just have $\Hyp_e(G)=\{[G]\}$. Thus when describing $\Hyp(G)$ for a given group, we just need to understand $\Hyp_\ell(G)$, $\Hyp_{qp}(G)$, and $\Hyp_{gt}(G)$.


\subsection{Pseudocharacters}\label{ssec:pseudochars}

In this subsection we recall some background on Busemann pseudocharacters of certain hyperbolic actions.

\begin{defi}[Quasimorphism, pseudocharacter, homogenization]
Let $G$ be a group and $q\colon G\to\R$ a function. We call $q$ a \emph{quasimorphism} if there exists a constant $D$ such that
\[
|q(gh)-q(g)-q(h)|\le D
\]
for all $g,h\in G$. We say that $q$ has \emph{defect} at most $D$. Note that a quasimorphism with defect $0$ is just a homomorphism to $\R$. If a quasimorphism becomes a homomorphism upon restriction to each cyclic subgroup, we call it a \emph{pseudocharacter}. Given any quasimorphism $q$, the function $p(g)\coloneqq \displaystyle \lim_{n\to\infty}\frac{q(g^n)}{n}$ is a pseudocharacter, called the \emph{homogenization} of $q$.
\end{defi}

\begin{rema}
We will only be interested in a specific pseudocharacter, called the Busemann pseudocharacter, defined below. This is one of the few cases where the terminology ``pseudocharacter'' has stuck; in most other contexts, the more common terminology is ``homogeneous quasimorphism.''
\end{rema}

The special case of defect $0$ is important:

\begin{defi}[Character, character sphere]
A quasimorphism with defect $0$, i.e., a homomorphism $G \to \R$, is called a \emph{character} of $G$. Two characters are \emph{equivalent} if they are positive scalar multiples of each other. The set of equivalence classes of non-trivial characters of $G$ forms the \emph{character sphere} $\Sigma(G)$ of $G$. A character is \emph{discrete} if its image is cyclic.
\end{defi}

Most of the groups we will care about here have the following property:

\begin{defi}
We will call a group \emph{agreeable} if every pseudocharacter is a character.
\end{defi}

In fact, a group is agreeable if and only if every element in the commutator subgroup has vanishing stable commutator length, thanks to Bavard duality~\cite{bavard91}. One of the most prominent examples of agreeable groups is the following:

\begin{cit}[{\cite[Proposition~2.65]{scl}}]\label{amenable:agreeable}
Amenable groups are agreeable. In particular, abelian groups are agreeable.
\end{cit}

A further example, which is the most relevant for our purposes, is groups of piecewise linear homeomorphisms of the interval~\cite{calegari:PL}; we will return to this point later on. The pseudocharacters we will be interested in -- the so called Busemann pseudocharacters -- arise from hyperbolic actions fixing a point in the boundary, as we now describe.

\begin{defi}[Busemann pseudocharacter]
Let $G$ be a group acting on a hyperbolic space $X$. Suppose the action fixes a point $\xi\in\partial X$ in the boundary. Let $(x_i)_{i\in\N}$ be a sequence of points in $X$ converging to $\xi$, and let $x\in X$. The function
\[
g\mapsto \displaystyle \limsup_{n\to\infty} \hspace{3pt} (d(x,x_n) - d(gx,x_n))
\]
is a quasimorphism, and we denote by $p\colon G\to\R$ its homogenization. It turns out that $p$ is independent of the choices of the $x_i$ and $x$, and only depends on $\xi$. We call $p$ the \emph{Busemann pseudocharacter} corresponding to this action of $G$ and the fixed point $\xi$.
\end{defi}

\begin{rema}\label{rem:left_right}
Our definition of the Busemann pseudocharacter here is the negative of the definition in~\cite{confining}. The reason is that, throughout the present paper we use right conjugation, whereas in~\cite{confining} they use left conjugation, so this is the convention that will make future citations to~\cite{confining} accurate. For the record, in~\cite{caprace15} the Busemann pseudocharacter is defined in the same way we define it here. Also note that with our convention, $\xi$ is an attracting (rather than repelling) fixed point for elements with positive $p$ value.
\end{rema}

The most important property of the Busemann pseudocharacter is the following:

\begin{cit}[{\cite[Section 4.1]{manning}}]
\label{cit:busemann:loxodromic}
    Let $G$ be a group acting on a hyperbolic space $X$ fixing $\xi \in \partial X$, and let $p$ be the Busemann pseudocharacter. Then $p(g) \neq 0$ if and only if $g$ is loxodromic. Therefore, $p$ is unbounded as soon as $G$ contains a loxodromic element.
\end{cit}

If the action is quasi-parabolic then there is a unique fixed point $\xi$, so we can simply refer to ``the Busemann pseudocharacter'' of the action. If the action is oriented lineal then there are two fixed points, but the resulting Busemann functions are negatives of each other, so again we will just say ``the Busemann pseudocharacter'', with an implicit sign ambiguity.


\subsection{Non-oriented lineal actions}

Let $G$ be a group with a non-oriented lineal action on a hyperbolic space $X$. Since the action on limit points is non-trivial, we cannot define the Busemann pseudocharacter. However, we can still define a meaningful object as follows.

\begin{defi}[Quasicocycle, defect, cocycle]
    Let $G$ be a group, $\varepsilon \colon G \to (\{ \pm 1 \}, \times)$ a homomorphism, and $\varphi \colon G \to \R$ a function. We call $\varphi$ an \emph{$\varepsilon$-quasicocycle} if there exists a constant $D$ such that
    \[ |\varphi(gh) - \varphi(g) - \varepsilon(g)\varphi(h)| \leq D\]
    for all $g, h \in G$. We say that $\varphi$ has \emph{defect} at most $D$. Note that $\varphi$ has defect $0$ if and only if the map $(\varphi, \varepsilon) \colon G \to \R \rtimes \{ \pm 1 \}$ is a homomorphism, in which case we call $\varphi$ an \emph{$\varepsilon$-cocycle}. Note also that if $\varepsilon$ is the trivial homomorphism, then we recover the definition of quasimorphism.
\end{defi}

\begin{cit}[{\cite[Proposition 2.10]{NL}}]
\label{cit:quasicocycle}
    Let $G$ be a group with a non-oriented lineal action on a hyperbolic space $X$ with limit points $\xi_{\pm}$. Let $\varepsilon(g)$ be the sign of the permutation that $G$ induces on $\{ \xi_+, \xi_- \}$. Let $K$ be the kernel of $\varepsilon$. Then the Busemann pseudocharacter $p \colon K \to \R$ extends to an $\varepsilon$-quasicocycle $\varphi \colon G \to \R$ of defect at most twice the defect of $p$. In particular $\varphi$ is an $\varepsilon$-cocycle if and only if $p$ is a character.
\end{cit}

Note that, while $\varepsilon$, $K$, and $p$ (up to sign) are canonically associated to the action, the construction of $\varphi$ requires a choice; this is due to the fact that the splitting $\mathrm{Isom}(\R) \cong \R \rtimes \{ \pm 1 \}$ is not canonical. We can now deduce the following criterion for excluding non-oriented lineal actions.

\begin{coro}\label{cor:nonoriented}
Let $G$ be a finitely generated agreeable group such that every index-$2$ subgroup of $G$ is agreeable. If $G$ does not surject onto $D_\infty$, then every lineal action of $G$ is oriented.
\end{coro}

\begin{proof}
Suppose that $G$ admits a non-oriented lineal action on the hyperbolic space $X$. Let $\varepsilon, K, p, \varphi$ be as in Citation~\ref{cit:quasicocycle}. Since $K$ is agreeable, $p$ is a character, and therefore $\varphi$ is an $\varepsilon$-cocycle, which gives us a homomorphism $(\varphi, \varepsilon) \ G \to \R \rtimes \{ \pm 1 \}$. Moreover, $K$ contains loxodromics, so $p$ is unbounded by Citation~\ref{cit:busemann:loxodromic}, and therefore $(\varphi, \varepsilon)$ has infinite image. Since $G$ is finitely generated, this produces a surjective homomorphism $G \to D_\infty$ (see the proof of~\cite[Proposition 3.1]{NL}).
\end{proof}

For studying actions of $F_n^{\pm}$, we will need the following more precise statement.

\begin{prop}
\label{prop:nonoriented:classification}

Let $G$ be a group, and let $X$ be a hyperbolic space with a non-oriented lineal action. Let $\varepsilon \colon G \to \{ \pm 1 \}$ and $p \colon \ker(\varepsilon) \to \R$ be as in Citation~\ref{cit:quasicocycle}. Then $\varepsilon$ and $p$ uniquely determine the equivalence class of $G \curvearrowright X$.
\end{prop}

\begin{proof}
By the classification of oriented lineal hyperbolic structures in terms of pseudocharacters~\cite[Section 2.3]{ABO}, the restricted action of $K \coloneqq \ker(\varepsilon)$ on $X$ is equivalent to the action on $\Cay(K, A)$, where
\[A \coloneqq \{ g \in K \mid |p(g)| \leq C \}\]
for $C > 0$ a large enough constant. By the Svarc--Milnor Lemma for cobounded actions~\cite[Lemma 3.11]{ABO}, the action $K \curvearrowright X$ corresponds to the generating set
\(A' = \{ g \in K \mid d_X(x, gx) \leq C' \},\)
where $C' > 0$ is another large enough constant, and $x$ is a point of $X$. Therefore the $A$-length of elements in $A'$ is uniformly bounded, and vice-versa.

Let $s \in G$ be such that $\varepsilon(s) = -1$. Now, using Svarc--Milnor again, the action $K \curvearrowright X$ corresponds to the generating set
\[B' = \{ g \in G \mid d_X(x, gx) \leq C' \};\]
up to taking a large enough constant $C'$, we can assume that $s \in B'$. Since every element of $G$ may be written uniquely as $ks^i$ for $k \in K, i \in \{0, 1\}$, we have
\[A' \cup \{ s \} \subseteq B' \subseteq A' \cup (A')^k s \]
for some large enough $k$. In other words the generating set $B'$ gives the same hyperbolic structure as the generating set $A' \cup \{ s \}$. Moreover, since $A$ and $A'$ are equivalent generating sets of $K$, also $A' \cup \{ s \}$ and $B \coloneqq A \cup \{ s \}$ are equivalent as generating sets of $G$. Since $B$ only depends on $\varepsilon$ and $p$, we conclude.
\end{proof}

\section{Confining subsets}\label{sec:confining}

In~\cite{caprace15}, Caprace, Cornulier, Monod, and Tessera introduced the notion of a confining subset in a group of the form $G=H\rtimes \Z$. This is a subset $Q\subseteq H$ such that the conjugation action of $\Z$ on $H$ satisfies certain properties with respect to $Q$, similar to the behavior of an ascending HNN-extension. In~\cite{confining}, the first author together with Abbott and Rasmussen generalized this to groups of the form $G=H\rtimes\Z^n$, now with a dependence on a choice of character $\rho\colon \Z^n\to\R$. These confining subsets for groups of the form $G=H\rtimes\Z^n$ turn out to be the key to producing lineal and quasi-parabolic actions of $G$. As a remark, here we will always view semidirect products as internal, so for example in $H\rtimes\Z^n$ the action of $\Z^n$ on $H$ is just the implicitly understood conjugation action. In~\cite{caprace15,confining}, semidirect products tend to be external, written $H\rtimes_\gamma\Z^n$ for some $\gamma\colon \Z^n\to\Aut(H)$. We will use the standard notation $x^y\coloneqq y^{-1}xy$ for conjugation; see Remark~\ref{rem:left_right}.

As a remark, we will always write semidirect products with the normal subgroup on the left, as is standard, though keep in mind that we tend to use right conjugation actions. Since we always view semidirect products internally, this will never matter.


\subsection{Definitions and properties}\label{ssec:confining_defs}

In this subsection we recall the definition of confining subset, and inspect some important properties.

\begin{defi}[Confining subset]\label{def:confining}
Let $G=H\rtimes \Z^n$ and let $\rho\colon \Z^n \to\R$ be a character. Call a subset $Q\subseteq H$ \emph{confining with respect to $\rho$}, or \emph{$\rho$-confining} if it is symmetric (i.e., $Q = Q^{-1}$) and the following properties hold:
\begin{description}
    \listlabel{StayInConf} For all $z\in\Z^n$ with $\rho(z)\ge 0$ we have $Q^z\subseteq Q$.
    \listlabel{GetInConf} For all $h\in H$ there exists $z\in \Z^n$ such that $h^z\in Q$.
    \listlabel{ProdConf} There exists $z_0\in \Z^n$ such that $(Q\cdot Q)^{z_0}\subseteq Q$.
\end{description}
We call $Q$ \emph{confining} if it is confining with respect to some $\rho$. We say $Q$ is a ``confining subset of'' either $H$ or $G$, depending on context.
\end{defi}

\begin{rema}[First consequences]
\label{rem:conf:consequences}
We note some direct consequences of the definition.
\begin{enumerate}
\item Since $Q$ is symmetric, \ProdConf\ implies that $1\in Q$.
\item If $Q$ is confining, then $G$ is generated by $Q$ together with $\Z^n$, thanks to \GetInConf.
\end{enumerate}
\end{rema}

Call $Q$ \emph{strictly confining} if $Q^z\subsetneq Q$ for some $z\in\Z^n$. Note that if $Q$ is confining but not strictly confining, then necessarily $Q=H$, and indeed $H$ is confining with respect to all $\rho$. If $Q$ is a confining subset that is a subgroup of $H$ we call it a \emph{confining subgroup}; in this case note that \ProdConf\ holds trivially.

\begin{rema}
    \emph{Confining subgroups} should not be confused with the better-studied notion of \emph{confined subgroup}~\cite{confined:subgroup}.
\end{rema}

A very useful fact in all that follows is that we actually have quite a bit of control over the conjugating elements in \GetInConf\ and \ProdConf. More precisely, the elements can be drawn from certain cyclic subgroups, as this next result shows.

\begin{lemm}
\label{lem:conf:cyclic}

Let $\rho \colon \Z^n \to \R$ be a character, and let $Q \subseteq H$ be a $\rho$-confining subset. Let $a \in \Z^n$ be any element such that $\rho(a) > 0$. Then:
\begin{enumerate}
\item For all $h \in H$ there exists $k \geq 0$ such that $h^{a^k} \in Q$;
\item There exists $k \geq 0$ such that $(Q \cdot Q)^{a^k} \subseteq Q$.
\end{enumerate}
\end{lemm}

\begin{proof}
For the first item, \GetInConf\ gives an element $z \in \Z^n$ such that $h^z \in Q$. Let $k \geq 0$ be large enough that $\rho(a^k) \geq \rho(z)$. Then $h^{a^k} = (h^z)^{a^k z^{-1}} \in Q^{a^k z^{-1}} \subseteq Q$, where in the last step we used \StayInConf. The second item is proved similarly.
\end{proof}

Another useful consequence of the definitions is the following strengthening of \ProdConf.

\begin{lemm}
\label{lem:products}

Let $\rho \colon \Z^n \to \R$ be a character, and let $Q \subseteq H$ be a $\rho$-confining subset. Let $z_0 \in \Z^n$ be given by \ProdConf, and suppose that $\rho(z_0) \geq 0$. Then $(Q^k)^{z_0^{k-1}} \subseteq Q$.
\end{lemm}

Note that by Lemma~\ref{lem:conf:cyclic}, this shows that the conjugating element taking $Q^k$ into $Q$ can always be chosen to be a power of a given $a \in \Z^n$ such that $\rho(a) > 0$.

\begin{proof}
We proceed by induction on $k$, the base case $k=1$ being trivially true. For $k\ge 2$ we have $(Q^k)^{z_0^{k-1}} = ((Q^{k-1})^{z_0^{k-2}} \cdot Q^{z_0^{k-2}})^{z_0}$. By induction $(Q^{k-1})^{z_0^{k-2}}\subseteq Q$, and by \StayInConf\ $Q^{z_0^{k-2}}\subseteq Q$. Thus, $(Q^k)^{z_0^{k-1}}\subseteq (Q \cdot Q)^{z_0}$, which is contained in $Q$ by \ProdConf.
\end{proof}

Next we would like a way to compare two confining subsets. Given a character $\rho\colon \Z^n\to\R$, we define a subset $Z_\rho$ of $\Z^n$ as follows. Fix a generating set $\{a_1,\dots,a_n\}$ for $\Z^n$, and a constant $C_\rho>0$ such that $|\rho(a_i)| \leq C_\rho$ for all $1\le i\le n$ and such that there exists $z_{top}\in \Z^n$ with $\rho(z_{top})=C_\rho$. Now set
\[
Z_\rho \coloneqq \{z\in\Z^n \mid |\rho(z)|\le C_\rho\}\text{,}
\]
so in particular $Z_\rho$ generates $\Z^n$. (Note that for $n\ge 2$ this is an infinite generating set.) As a remark, the Cayley graph of $\Z^n$ with respect to the generating set $Z_\rho$ is quasi-isometric to a line, and the action of $\Z^n$ on this hyperbolic space is oriented lineal. We also have that for $G=H\rtimes \Z^n$, if $Q\subseteq H$ is $\rho$-confining for some $\rho\colon \Z^n\to\R$, then $Q\cup Z_\rho$ is a symmetric generating set for $G$. Choosing a different finite generating set for $\Z^n$ and/or a different constant $C_\rho$ can change the generating set $Q\cup Z_\rho$ for $G$, but cannot change it up to the equivalence relation on generating sets of $G$, so up to equivalence we are always free to choose the $a_i$ and $C_\rho$ as we see fit. Also note that we can always assume $Z_{-\rho}=Z_\rho$.

Consider the following relation on confining subsets. Given $G=H\rtimes \Z^n$ and $\rho\colon\Z^n\to\R$ as before, if $Q_1,Q_2\subseteq H$ are $\rho$-confining subsets, we write
\[
Q_1\preceq_\rho Q_2
\]
if $Q_2$ is bounded in the alphabet $Q_1\cup Z_\rho$. Note that $Q_1\preceq_\rho Q_2$ for $\rho$-confining subsets is clearly equivalent to $Q_1\cup Z_\rho\preceq Q_2\cup Z_\rho$ for generating sets. Since $\preceq_\rho$ is a preorder, we get an equivalence relation $\sim_\rho$ on $\rho$-confining subsets and an induced partial order $\preccurlyeq_\rho$ on equivalence classes $[Q]$ of $\rho$-confining subsets $Q$.

\begin{defi}[The posets $\Conf_\rho(G)$ and $\Conf^\ast_\rho(G)$]\label{def:conf_poset}
The poset of equivalence classes of $\rho$-confining subsets with the partial order $\preccurlyeq_\rho$ is denoted by $\Conf_\rho(G)$. The subposet of \emph{strictly} $\rho$-confining subsets is denoted by $\Conf^\ast_\rho(G)$. Note that $\Conf_\rho(G)$ is a cone over $\Conf^\ast_\rho(G)$, with cone point the equivalence class of $H$.
\end{defi}

If $Q_1 \subseteq Q_2$ then clearly $Q_2 \preceq_\rho Q_1$. More generally, it turns out that if $Q_1^z \subseteq Q_2$ for some $z\in \Z^n$ then $Q_2 \preceq_\rho Q_1$, and, perhaps surprisingly, the converse of this is true. This will be an easy consequence of the following characterization of boundedness in the alphabet $Q\cup Z_\rho$ for confining $Q$.

\begin{prop}\label{prop:conj_is_all}
Let $Q\subseteq H$ be a $\rho$-confining subset. Then a subset $B\subseteq H$ is bounded in the alphabet $Q\cup Z_\rho$ if and only if $B \subseteq Q^z$ for some $z \in \Z^n$.
\end{prop}

\begin{proof}
First suppose $B \subseteq Q^z$. Let $m$ be the length of $z$ in the alphabet $Q\cup Z_\rho$. Now observe that every element of $B$ can be written as $z^{-1}qz$ for some $q\in Q$, and this has length $2m+1$ in the alphabet $Q\cup Z_\rho$.

Now suppose that $B$ is bounded in the alphabet $Q\cup Z_\rho$. Say every element of $B$ can be written as a product of at most $m$ elements of $Q\cup Z_\rho$. In order to only have to deal with a single form, note that this implies every element $b$ of $B$ can be written in the form $b=q_1 z_1\cdots q_m z_m$ for some $q_i\in Q$, $z_i\in Z_\rho$ (with several $q_i$ and/or $z_i$ possibly being the identity).
Now fix $z_{top}\in Z_\rho$ with $\rho(z_{top})=C_\rho$, as in the construction of $Z_\rho$. For any $z\in Z_\rho$, since $|\rho(z)|\le C_\rho$ we have $\rho(z_{top} z)\ge 0$. In particular, $Q^{z_{top} z}\subseteq Q$ for all $z\in Z_\rho$ by \StayInConf, so $Q^z\subseteq Q^{z_{top}^{-1}}$.
More generally, for any $k \in \Z$ we have $(Q^{z_{top}^{-k}})^z\subseteq Q^{z_{top}^{-k-1}}$ for all $z\in Z_\rho$, as can be seen by conjugating both sides of this expression by $z_{top}^{k+1}$ by \StayInConf.


Now we return to our arbitrary $b\in B$ with $b=q_1 z_1\cdots q_m z_m$. Since each $z_i$ is in $Z_\rho$, the above allows us to sequentially move all the $z_i$ to the left, yielding $b=z_1\cdots z_m q_1'\cdots q_m'$, with $q_m'\in Q^{z_{top}^{-1}}$, $q_{m-1}'\in Q^{z_{top}^{-2}}$, and so forth up to $q_1'\in Q^{z_{top}^{-m}}$. Since $b\in H$, we must have $z_1\cdots z_m=1$, and since $Q^{z_{top}^{-k}}\subseteq Q^{z_{top}^{-m}}$ for all $0\le k\le m$ by \StayInConf, we get $b\in (Q^m)^{z_{top}^{-m}}$. Since $b$ was arbitrary, this implies that $B\subseteq (Q^m)^{z_{top}^{-m}}$. By Lemma~\ref{lem:products} we can choose $z' \in \Z^n$ such that $(Q^m)^{z'}\subseteq Q$, and so we conclude that $B\subseteq Q^{z^{-1}}$ where $z =  z_{top}^mz' \in \Z^n$.
\end{proof}

\begin{coro}[Characterization of $\preceq_\rho$]\label{cor:conjugate}
Let $Q_1,Q_2\subseteq H$ be $\rho$-confining subsets, and fix $a \in \Z^n$ such that $\rho(a) > 0$. Then $Q_1 \preceq_\rho Q_2$ if and only if $Q_2^{a^k} \subseteq Q_1$ for some~$k \geq 0$.
\end{coro}

\begin{proof}
By Proposition~\ref{prop:conj_is_all}, we know that $Q_1 \preceq_\rho Q_2$ if and only if $Q_2^z \subseteq Q_1$ for some $z \in \Z^n$. This gives one direction; for the other one, let $k \geq 0$ be large enough that $\rho(a^k) \geq \rho(z)$. Then
\[Q_2^{a^k} = (Q_2^z)^{a^k z^{-1}} \subseteq Q_1^{a^k z^{-1}} \subseteq Q_1,\]
thanks to \StayInConf, since $\rho(a^k z^{-1}) \geq 0$.
\end{proof}

To the best of our knowledge, Corollary~\ref{cor:conjugate} has not appeared before in the literature, but it encapsulates a very natural and straightforward characterization of the equivalence relation on confining subsets, and the partial order on their equivalence classes. In particular, this characterization facilitates the proof of a structural property of the poset $\Conf^\ast_\rho(G)$. Recall that the \emph{join} of two elements in a poset is their least upper bound, and a poset is a \emph{join semilattice} if every two elements admit a (necessarily unique) join.

\begin{prop}
\label{prop:join}

Let $G=H\rtimes \Z^n$ and let $\rho\colon \Z^n \to\R$ be a character. The poset $\Conf^\ast_\rho(G)$ is a join semilattice. Given $\rho$-confining subsets $Q_1, Q_2 \subseteq H$, the join of $[Q_1]$ and $[Q_2]$ is $[Q_1 \cap Q_2]$.
\end{prop}

\begin{proof}
First we show that $Q_1 \cap Q_2$ is also a $\rho$-confining subset. For $z \in \Z^n$ with $\rho(z) \geq 0$ we have $(Q_1 \cap Q_2)^z = Q_1^z \cap Q_2^z \subseteq Q_1 \cap Q_2$, which proves \StayInConf.
Let $h \in H$ and let $z_i$ ($i=1,2$) be such that $h^{z_i} \in Q_i$. By Lemma~\ref{lem:conf:cyclic} we may choose $z \in \Z^n$ that works as both $z_1$ and $z_2$. Then $h^z \in Q_1 \cap Q_2$, which proves \GetInConf.
Finally, let $z_i$ ($i=1,2$) be such that $(Q_i \cdot Q_i)^{z_i} \in Q_i$. Again, we may choose $z \in \Z^n$ that works as both $z_1$ and $z_2$. Then
\[
((Q_1 \cap Q_2) \cdot (Q_1 \cap Q_2))^{z} \subseteq ((Q_1 \cdot Q_1) \cap (Q_2 \cdot Q_2))^{z} \subseteq Q_1 \cap Q_2,
\]
which proves \ProdConf.

Clearly $[Q_1 \cap Q_2]$ is a common upper bound of $[Q_1]$ and $[Q_2]$. Let $R$ be a $\rho$-confining subset such that $[R]$ is another common upper bound. By Corollary~\ref{cor:conjugate} there exist $z_1, z_2 \in \Z^n$ such that $R \subseteq Q_1^{z_1} \cap Q_2^{z_2}$. Suppose without loss of generality that $\rho(z_1) \geq \rho(z_2)$. Then $\rho(z_1 z_2^{-1}) \geq 0$ and so \[Q_1^{z_1} = (Q_1^{z_1 z_2^{-1}})^{z_2} \subseteq Q_1^{z_2}.\]
This shows that $R \subseteq (Q_1 \cap Q_2)^{z_2}$ and so by Corollary~\ref{cor:conjugate} again, $[Q_1 \cap Q_2] \preccurlyeq_\rho [R]$.
\end{proof}

\begin{rema}
In some cases that have been considered before, $\Conf^\ast_\rho(G)$ is also a meet semi-lattice, and thus a lattice~\cite{sahana:wreath, AR:BS, ABR2}. However this is not true in general, and we will see that the groups $F_n$ provide an example (Remark~\ref{rem:meet}).
\end{rema}

In practice, the confining subsets one tends to encounter in the literature are always equivalent to a confining subgroup, see for example~\cite{caprace15, sahana:wreath, AR:BS, ABR2}. This can be characterized as follows:

\begin{prop}\label{prop:equiv_to_subgroup}
Let $G=H\rtimes \Z^n$ and let $\rho\colon \Z^n \to\R$ be a character. Let $Q\subseteq H$ be a $\rho$-confining subset. Then the following are equivalent:
\begin{enumerate}
    \item $Q$ is equivalent to a $\rho$-confining subgroup.
    \item $Q$ is equivalent to $\langle Q\rangle$.
    \item There exists $z\in \Z^n$ such that $\langle Q\rangle^z\subseteq Q$.
\end{enumerate}
\end{prop}

\begin{proof}
It follows easily from the definitions that if $Q$ is a confining subset, then $\langle Q \rangle$ is a confining subgroup. Clearly (ii) implies (i), and Corollary~\ref{cor:conjugate} shows that (ii) and (iii) are equivalent, so it suffices to prove that (i) implies (iii). If $Q$ is equivalent to a $\rho$-confining subgroup $R$, then by Corollary~\ref{cor:conjugate} we have $Q^z\subseteq R$ and $R^{z'}\subseteq Q$ for some $z,z'\in\Z^n$. Since $R$ is a subgroup we have $\langle Q\rangle^z\le R$, so $\langle Q\rangle^{zz'}\subseteq Q$ and we are done.
\end{proof}

We will see however that lamplighters with infinite base group admit (many) confining subsets that are not confining subgroups (Example~\ref{ex:balls}, Example~\ref{ex:machado}), and we will see more examples of this among non-lamplike confining subsets of $F_n$ (Section~\ref{sec:moving}).


\subsection{Confining subsets and hyperbolic actions}\label{ssec:conf_hyp}

We first make the following observation about pseudocharacters of groups of the form $H \rtimes \Z^n$:

\begin{lemm}
\label{lem:factor}
Let $p \colon H \rtimes \Z^n \to \R$ be a pseudocharacter. If $p(H) = 0$, then $p$ factors as $H \rtimes \Z^n \to \Z^n \xrightarrow{\rho} \R$ for some character $\rho$ of $\Z^n$.
\end{lemm}

\begin{proof}
Since $p(H) = 0$, the left exactness of the pseudocharacter functor~\cite[Remark 2.90]{scl} implies that $p$ factors through a pseudocharacter $\rho \colon \Z^n \to \R$. But $\Z^n$ is agreeable (Citation~\ref{amenable:agreeable}), so $\rho$ is actually a character.
\end{proof}

Now we spell out the connection between confining subsets and hyperbolic actions.

\begin{cit}[Confining subsets to hyperbolic actions {\cite[Theorem~1.1]{confining}}]\label{cit:confining_to_action}
Let $G=H\rtimes \Z^n$, and let $\rho\colon \Z^n\to \R$ be a non-trivial character. If $Q\subseteq H$ is confining with respect to $\rho$, then $\Cay(G,Q\cup Z_\rho)$ is hyperbolic. Moreover, if $Q$ is strictly confining then $[Q\cup Z_\rho]\in \Hyp_{qp}(G)$, and if $Q$ is not strictly confining then $[Q\cup Z_\rho]\in \Hyp_\ell^+(G)$. Finally, the Busemann pseudocharacter for $G\curvearrowright \Cay(G,Q\cup Z_\rho)$ sends $H$ to $0$, so induces a character $\Z^n\to\R$, which is equivalent to $\rho$ in the quasi-parabolic case and to $\pm\rho$ in the oriented lineal case.
\end{cit}

In~\cite[Theorem~1.1]{confining} the last claim in the quasi-parabolic case is that the character is ``proportional'' to $\rho$, but the proof of~\cite[Lemma~1.1]{confining} makes clear that it is in fact equivalent to $\rho$, meaning the constant of proportionality is positive. Also note that the Busemann pseudocharacter for $G\curvearrowright \Cay(G,Q\cup Z_\rho)$ is in fact a character.

\begin{cit}[Hyperbolic actions to confining subsets {\cite[Theorem~3.17]{confining}}]
\label{cit:action_to_confining}
Let $G=H\rtimes \Z^n$, and suppose $G$ acts isometrically on a hyperbolic space $X$. Assume the action is cobounded, and either oriented lineal or quasi-parabolic. Let $p\colon G\to\R$ be the Busemann pseudocharacter associated to this action, and assume that $p(H)=0$, so $p$ induces a (non-trivial) character $\rho\colon \Z^n\to\R$. Then there exists a subset $Q\subseteq H$ that is confining with respect to $\rho$, such that $X$ is $G$-equivariantly quasi-isometric to the Cayley graph $\Cay(G,Q\cup Z_\rho)$. If the action is oriented lineal then $Q=H$, and if the action is quasi-parabolic then $Q$ is strictly confining.
\end{cit}

Note that we have specified that the characters $\rho$ used above are non-trivial. If $\rho=0$ then $Z_0=\Z^n$, and $Q=H$ is the only confining subset with respect to $\rho=0$, so $[Q\cup Z_0]=[Q\cup \Z^n]=[G]$. In other words, $\rho=0$ corresponds precisely to $\Hyp_e(G)=\{[G]\}$, and this is the end of the story for the case when $\rho$ is trivial.

Before we begin imposing conditions on our groups, let us prove two more general results, which are specializations of Citation~\ref{cit:confining_to_action} and Citation~\ref{cit:action_to_confining} that establish a relationship between confining subgroups and actions on trees (\emph{all trees are assumed to be simplicial}). For the next two proofs, we assume some familiarity with Bass--Serre theory~\cite{serre}.

\begin{prop}[Confining subgroups yield trees]\label{prop:subgroups_to_trees}
Let $G=H\rtimes \Z^n$, and let $\rho\colon \Z^n\to \R$ be a (non-trivial) discrete character. Let $Q$ be a $\rho$-confining subset that is a subgroup. Then the hyperbolic structure $[Q\cup Z_\rho]$ arising from Citation~\ref{cit:confining_to_action} is represented by a simplicial action on a tree.

More precisely, it is represented by the action on the Bass--Serre tree of an ascending HNN-extension expression of $G$, with base group $Q \ker(\rho)$, and stable letter an element $a \in \Z^n$ such that $\rho(a)$ is positive and minimal.
\end{prop}

\begin{proof}
Since the character is discrete, we may assume that it is a surjection $\rho \colon \Z^n \to \Z$. Choose a basis $\{a_0,\dots,a_{n-1}\}$ for $\Z^n$ such that $\rho(a_0)=1$ and $\rho(a_i)=0$ for all $1\le i\le n-1$. We are free to choose $C_\rho=1$, so that $Z_\rho=\{a_0^{k_0}\cdots a_{n-1}^{k_{n-1}}\mid k_0\in\{-1,0,1\}\}$. Let $K\coloneqq Q\langle a_1,\dots,a_{n-1}\rangle$; since $\rho(a_i)=0$ for all $1\le i\le n-1$, each such $a_i$ normalizes $Q$ by \StayInConf, and since we are assuming that $Q$ is a subgroup of $G$, we get that $K$ is a subgroup of $G$.

We claim that $G$ decomposes as an ascending HNN-extension with base $K$ and stable letter $a_0$. For this we need to check the three conditions from~\cite[Lemma~3.1]{ascending}. First, $K^{a_0}\le K$ since $Q^{a_0}\le Q$ by \StayInConf. Second, $G$ is generated by $K\cup\{a_0\}$, since $Q\le K$, $Z_\rho\subseteq \{1,a_0,a_0^{-1}\}K$, and $G$ is generated by $Q\cup Z_\rho$. Finally, $\langle a_0\rangle\cap K=\{1\}$ since $\rho(a_0)=1$ and $K$ lies in the kernel of the composition of $\rho$ with the projection $G\to\Z^n$. This concludes the proof of the claim.

Let $X$ be the Bass--Serre tree of this ascending HNN-extension expression of $G$, and let $v \in V(X)$ be the vertex whose stabilizer is $K$. The action of $G$ on $X$ has a fundamental domain of diameter $1$, so by the Svarc--Milnor Lemma for cobounded actions~\cite[Lemma 3.11]{ABO}, this action corresponds to the hyperbolic structure with generating set
\(B = \{ g \in G \mid d(v, gv) \leq 3 \}.\)
We note that $d_T(v, gv) \leq 1$ for every $g \in A = Q \cup Z_\rho$, so $A \subseteq B$. By Bass--Serre theory, every element of $B$ can be written as a product of powers of $a_0$ and elements in $K$, and the length of this product is bounded uniformly because the distance from $v$ to $gv$ for $g\in B$ is bounded uniformly by $3$. It follows that $A$ and $B$ are generating sets representing the same hyperbolic structure, which is the one coming from the action of $G$ on this Bass--Serre tree.
\end{proof}

\begin{prop}[Trees yield confining subgroups] \label{prop:trees_to_subgroups}
Let $G=H\rtimes \Z^n$, and suppose $G$ acts simplicially on a tree $X$. Assume the action is cobounded, and either oriented lineal or quasi-parabolic. Assume the Busemann pseudocharacter associated to this action sends $H$ to $0$, and so induces a (non-trivial) character $\rho\colon \Z^n\to\R$. Then $\rho$ is discrete and the $\rho$-confining subset arising from Citation~\ref{cit:action_to_confining} can be chosen to be a subgroup.
\end{prop}

\begin{proof}
Since the Busemann pseudocharacter sends $H$ to $0$, every element of $H$ is elliptic, as trees do not admit parabolic isometries. The restriction to the $\Z^n$ retract is an action of an abelian group on a tree that admits a loxodromic element, so it must be oriented lineal (see e.g.~\cite[Proposition 3.2]{NL}). Since the action is on a tree, there exists a $\Z^n$-invariant line, which we denote by $L$, and the image of $\rho$ is equal to $\langle \ell \rangle$, where $\ell \in \N$ is the minimal positive translation length of an element of $\Z^n$. Because the action is simplicial, $\rho$ is discrete. 
If the action is oriented lineal then our confining subset is $H$, which is a subgroup and so we are done.

Now assume the action is quasi-parabolic. Let $\xi$ be the unique $G$-fixed point in the boundary of $X$, and note that the $\Z^n$-invariant line $L$ has $\xi$ as one of its endpoints. By the definition of the Busemann pseudocharacter, if $z \in \Z^n$ satisfies $\rho(z) > 0$, then $z^n x \xrightarrow{n \to \infty} \xi$ for all $x \in X$ (Remark~\ref{rem:left_right}). Choose $z_0\in\Z^n$ such that $\rho(z_0) = \ell > 0$. Since the action of $G$ on $X$ is cobounded, we can choose a bounded fundamental domain for the action, say with vertex set $D$, and without loss of generality we can assume that $v_0$ lies in $L$ for some $v_0\in D$. For any $v$ in $X$, write $[v,\xi)$ for the unique ray from $v$ to $\xi$.

Since $D$ is bounded, the geodesic projection of $D$ onto $L$ is also bounded, and so we can choose a vertex $w$ of $L$ such that $w$ lies in $[v,\xi)$ for all $v\in D$. Since $\xi$ is fixed by all of $G$, we see that $w$ is fixed by $\Stab_G(v)$ for all $v\in D$. Set $Q=\Stab_H(w)$, and we claim that $Q$ is a $\rho$-confining subgroup. Being a subgroup, $Q$ is symmetric and satisfies \ProdConf. To verify \StayInConf, note that for $z\in \Z^n$ with $\rho(z)\ge 0$ and $q\in Q$, we have that $zw$ lies in $[w,\xi)$, and so $q(zw)=zw$, thus confirming that $z^{-1}qz\in Q$. For \GetInConf, note that for any $h\in H$ the lines $L$ and $h(L)$ must contain a common ray to $\xi$, which is fixed pointwise by $h$ because $\rho(h) = 0$; in particular $h$ fixes some vertex $w'$ of $L$ lying in $[w,\xi)$. Up to choosing $w'$ closer to $\xi$, there exists $k \geq 0$ such that $z_0^k w = w'$, and hence $z_0^{-k}hz_0^k\in Q$.

It remains to prove that $Q$ corresponds to $G\curvearrowright X$ under the correspondence from Citation~\ref{cit:action_to_confining}. For this, we need to show that the generating set $Q \cup Z_\rho$, where $Z_\rho = \{ z \in \Z^n \mid |\rho(z)| \leq \ell \}$, is equivalent to the generating set given by the Svarc--Milnor Lemma for cobounded actions~\cite[Lemma 3.11]{ABO}, namely $B \coloneqq \{ g \in G \mid d(v_0, gv_0) \leq 3d \}$, where $d$ is the diameter of our bounded fundamental domain with vertex set $D$. Since all values of $\rho(G)$ are attained by $\langle z_0 \rangle$, by Bass--Serre theory the generating set $B$ is equivalent to the generating set
\[A \coloneqq \{ z_0^{\pm 1} \} \cup \bigcup_{v\in D} \Stab_G(v).\]
Next, note that $\Stab_G(w) = \Stab_H(w) \Stab_{\Z^n}(w) = Q \ker(\rho)$. Since $\ker(\rho) \subseteq Z_\rho$, we see that the generating set $Q \cup Z_\rho$ is equivalent to the generating set
\[A' \coloneqq \{z_0^{\pm 1}\}\cup \Stab_G(w).\]

We are left to prove that $[A] = [A']$. On the one hand, $\Stab_G(v) \leq \Stab_G(w)$ for all $v\in D$, and so $A \subseteq A'$. In the other direction, choose $k \geq 0$ such that $w' \coloneqq z_0^k v_0$ lies in $[w,\xi)$; note that this is possible since $v_0$ was assumed to be in $L$. Then
\[z_0^{-k} \Stab_G(w) z_0^k \subseteq z_0^{-k} \Stab_G(w') z_0^k = \Stab_G(v_0),\]
which shows that $\Stab_G(w) \subseteq A^{2k+1}$, and therefore $A' \subseteq A^{2k+1}$; this concludes the proof.
\end{proof}

Thus, for discrete characters, the correspondence between confining subsets and hyperbolic actions restricts to a correspondence between confining subgroups and actions on trees.

\begin{rema}
\label{rem:simplicial}
    In Proposition~\ref{prop:trees_to_subgroups}, we assumed that the action on the tree $X$ is simplicial, and not merely by isometries. This is necessary to ensure that the character $\rho$ is discrete: a counterexample is the action defined via a dense embedding of $\Z^2$ into~$\R$, with its standard simplicial structure.

    However, this assumption is void in the case of quasi-parabolic actions, which is what we will be most interested in. Indeed, suppose that $G$ admits a quasi-parabolic action by isometries on a simplicial tree $X$ and let $g \in G$. Choose a vertex $v$ of valency at least $3$, so $gv$ must have the same valency and therefore be a vertex. Let $w$ be any other vertex, so $d(gv, gw) = d(v, w)$ is an integer, and hence $gw$ is also a vertex. This shows that every quasi-parabolic action by isometries on a simplicial tree is automatically simplicial.

    In particular, the second conclusion of Proposition~\ref{prop:trees_to_subgroups} holds for all actions by isometries on simplicial trees. If the action is oriented lineal, then the confining subset arising from Citation~\ref{cit:action_to_confining} is the whole of $H$, and if the action is quasi-parabolic, then it is automatically simplicial by the above, so we can apply Proposition~\ref{prop:trees_to_subgroups} as stated.
\end{rema}


\subsection{Oriented lineal and quasi-parabolic structures for agreeable, \texorpdfstring{$\Z^n$}{Zn}-split groups}\label{ssec:lin_qp}

The above results allow us to give a more explicit description of lineal and quasi-parabolic structures in the following cases:

\begin{defi}[$\Z^n$-split]
A group $G$ is \emph{$\Z^n$-split} if it admits a surjective map $G\to\Z^n$ that splits, so that $G=H\rtimes \Z^n$ for $H$ the kernel of this map. Call an expression $G=H\rtimes \Z^n$ a \emph{$\Z^n$-splitting} of a $\Z^n$-split group.
\end{defi}

Different choices of sections $\Z^n\to G$ will yield different $\Z^n$-splittings, which is why we are being careful to say ``a'' $\Z^n$-splitting. A special type of $\Z^n$-splitting will be most relevant for us:

\begin{defi}[Sharply $\Z^n$-split]
    A $\Z^n$-split $G = H \rtimes \Z^n$ is called \emph{sharp} if $H$ equals the kernel of the torsion-free abelianization of $G$. A group is \emph{sharply $\Z^n$-split} if it admits a sharp $\Z^n$-splitting.
\end{defi}

The relevance of sharp $\Z^n$-splittings lies in the following, which is just a consequence of the definitions.

\begin{lemm}\label{lem:sharp}
Let $G = H \rtimes \Z^n$ be a sharp $\Z^n$-splitting. Then every character of $G$ vanishes on $H$. If moreover $G$ is agreeable, then every pseudocharacter of $G$ vanishes on $H$. \qed
\end{lemm}

In particular, the hypothesis of Citation~\ref{cit:action_to_confining} is always satisfied for sharp $\Z^n$-splittings.

Our next goal is to characterize lineal and quasi-parabolic actions of agreeable, sharply $\Z^n$-split groups. First let us discuss the (much easier) lineal case.

\begin{prop}\label{prop:lineal}
Let $G$ be an agreeable, sharply $\Z^n$-split group. Let $G=H\rtimes\Z^n$ be a sharp $\Z^n$-splitting. Then the poset $\Hyp_\ell^+(G)$ is an antichain, and the function $\varphi \colon \Sigma(\Z^n)\to \Hyp_\ell^+(G)$ sending $[\rho]$ to $[H\cup Z_\rho]$ is a well defined 2-to-1 surjection with each fiber of the form $\{[\rho],[-\rho]\}$.
\end{prop}

\begin{proof}
The proof is inspired by the work in Subsection~4.2 of~\cite{confining}, which covers a special case. First note that $\Hyp_\ell^+(G)$ is an antichain simply because this is always the case~\cite[Theorem~4.22(a)]{ABO}. Now we claim $\varphi$ is well defined. This is because if $\rho$ and $\rho'$ are equivalent, then $H\cup Z_\rho$ and $H\cup Z_{\rho'}$ dominate each other, and the output $[H\cup Z_\rho]$ really lies in $\Hyp_\ell^+(G)$ by Citation~\ref{cit:confining_to_action}. Since every pseudocharacter of $G$ vanishes on $H$ by Lemma~\ref{lem:sharp}, Citation~\ref{cit:action_to_confining} applies, and gives us surjectivity. Finally, we need to show that $\varphi$ is 2-to-1 with the claimed fibers. Recall that we always assume $Z_{\rho} = Z_{-\rho}$, so $\varphi([\rho])=\varphi([-\rho])$. Now suppose $[\rho]\ne[\pm\rho']$, and we want to prove that $[H\cup Z_\rho]\ne[H\cup Z_{\rho'}]$. Since $[\rho]\ne[\pm\rho']$ we have $\ker(\rho)\ne \ker(\rho')$. Without loss of generality let us say that $\ker(\rho')\not\le \ker(\rho)$, and choose some $z\in \ker(\rho')\setminus \ker(\rho)$. Then $\langle z\rangle$ is a subgroup of $\ker(\rho')$, hence a subset of $Z_{\rho'}$, that is unbounded with respect to the word length coming from $H\cup Z_\rho$. We conclude that $H\cup Z_\rho\not\preceq H\cup Z_{\rho'}$, so $[H\cup Z_\rho]\ne[H\cup Z_{\rho'}]$.
\end{proof}

Now we focus on the quasi-parabolic case.

If $G = H \rtimes \Z^n$, and $\rho\colon\Z^n\to\R$ is a (non-trivial) character of $\Z^n$, write $\Hyp_{qp}^\rho(G)$ for the subposet of equivalence classes of those quasi-parabolic actions whose Busemann pseudocharacter induces a character $\Z^n\to\R$ equivalent to $\rho$.

\begin{prop}
\label{prop:qp}

Let $G$ be an agreeable, sharply $\Z^n$-split group. Let $\rho \colon \Z^n \to \R$ be a non-trivial character. Then the posets $\Hyp_{qp}^\rho(G)$ and $\Conf^\ast_\rho(G)$ are isomorphic.
\end{prop}

\begin{proof}
We claim that $[Q]\mapsto [Q\cup Z_\rho]$ is a well defined poset isomorphism from $\Conf^\ast_\rho(G)$ to $\Hyp_{qp}^\rho(G)$. Since $Q_1\preceq_\rho Q_2$ implies $Q_1\cup Z_\rho\preceq Q_2\cup Z_\rho$, it is a well defined poset map. Since $Q_1\cup Z_\rho\preceq Q_2\cup Z_\rho$ implies $Q_1\preceq_\rho Q_2$, it is injective. It just remains to prove that this map is surjective. Since $G$ is agreeable, Citation~\ref{cit:action_to_confining} says that every element of $\Hyp_{qp}^\rho(G)$ is of the form $[Q\cup Z_\rho]$ for some strictly confining $Q$, so we are done.
\end{proof}

Now we can tie together $\Hyp_\ell^+(G)$ and $\Hyp_{qp}^\rho(G)$.

\begin{prop}\label{prop:linvsqp}
Let $G$ be an agreeable, sharply $\Z^n$-split group. Let $\rho \colon \Z^n \to \R$ be a non-trivial character. The element $[H\cup Z_\rho]$ of $\Hyp_\ell^+(G)$ is a lower bound of every element of $\Hyp_{qp}^\rho(G)$, and is not related to any elements of any $\Hyp_{qp}^{\rho'}(G)$ for $[\rho']\ne[\pm\rho]$.
\end{prop}

\begin{proof}
The first claim is immediate since $[H\cup Z_\rho]\preccurlyeq [Q\cup Z_\rho]$ for all $\rho$-confining $Q$. We need to show the second claim. Suppose $[\rho']\ne [\pm\rho]$, so $\ker(\rho)\ne\ker(\rho')$. Extending $\rho$ and $\rho'$ from non-zero homomorphisms $\Z^n\to\R$ to non-zero linear transformations $\R^n\to\R$, the kernels $\ker(\rho)$ and $\ker(\rho')$ become distinct codimension-$1$ subspaces of $\R^n$. We can therefore find a sequence of points in $\ker(\rho')$ that gets arbitrarily far away from $\ker(\rho)$. Taking the $C_{\rho'}$-neighborhood and intersecting with $\Z^n$, we can choose a sequence of points in $Z_{\rho'}$ that gets arbitrarily far away from $Z_\rho$. This shows that for any $\rho'$-confining subset $Q$, the set $Q\cup Z_{\rho'}$ is unbounded in the alphabet $H\cup Z_{\rho}$, and so $[H\cup Z_\rho]\not\preccurlyeq[Q\cup Z_{\rho'}]$. Finally, it is obvious that $[Q\cup Z_{\rho'}]\not\preccurlyeq[H\cup Z_\rho]$ since Proposition~\ref{prop:lineal} says that $[H\cup Z_{\rho'}]\not\preccurlyeq [H\cup Z_\rho]$. We conclude that $[H\cup Z_\rho]$ is not related to any elements of any $\Hyp_{qp}^{\rho'}(G)$.
\end{proof}

Finally, let us prove one useful criterion that can be used to show that $\Conf^\ast_\rho(G)$ is empty, for certain $\Z^n$-splittings and certain $\rho$. Here a group is \emph{boundedly generated} by a collection of subgroups if it is generated by their union $U$, and is bounded with respect to the word length coming from the generating set $U$. For an element $g \in G$ we denote by $C_H(g) \coloneqq \{ h \in H \mid [g, h] = 1 \}$ its centralizer in $H$.

\begin{prop}
\label{prop:empty_conf}
Let $G = H \rtimes \Z^n$ be a $\Z^n$-splitting. Let $\{a_1,\dots,a_n\}$ be a basis for $\Z^n$ and $\rho\colon \Z^n \to \R$ a character. If $\rho(a_i) \neq 0$, then every $\rho$-confining subset contains $C_H(a_i)$. In particular, if $H$ is boundedly generated by those $C_H(a_i)$ with $\rho(a_i)\ne 0$, then $\Conf^\ast_\rho(G)= \emptyset$.
\end{prop}

\begin{proof}
Suppose that $\rho(a_i) \neq 0$, and let $c\in C_H(a_i)$. By Lemma~\ref{lem:conf:cyclic} there exists $k \in \Z$ such that $c^{a_i^k} \in Q$. But $c^{a_i^k} = c$ and so $C_H(a_i) \subseteq Q$.

If $H$ is boundedly generated by such $C_H(a_i)$, then $H=Q^\ell$ for some $\ell\in\N$. Now Lemma~\ref{lem:products} tells us that $Q$ is equivalent to $H$.
\end{proof}

\section{Brown--Higman--Thompson groups}\label{sec:thompson}

In this section we recall some background on the Brown--Higman--Thompson groups~$F_n$. We recall the following standard dynamical definitions, which we will repeatedly use in all that follows without further mention. For simplicity, we state them in our setting of one-dimensional actions:

\begin{defi}
\label{def:dynamical}

Let $G \leq \Homeo([0, 1])$. We will always view homeomorphisms as acting on $[0,1]$ from the right, so $x.g$ denotes the image of $x\in [0,1]$ under $g\in\Homeo([0,1])$.
\begin{enumerate}
\item An element $g \in G$ is \emph{orientation preserving} if $x < y$ implies $x.g < y.g$. If every element of $G$ is orientation preserving, we say that $G$ is \emph{orientation preserving}, and write $G \leq \Homeo^+([0, 1])$.
\item For an element $g \in G$, its \emph{support} is the set $\Supp(g) = \{ x \in [0, 1] \mid x.g \neq x \}$. For a set $S \subseteq G$, we denote $\Supp(S) = \cup_{g \in S} \Supp(g)$. We say that $S$ is \emph{supported on $I \subseteq [0, 1]$} if $\Supp(S) \subseteq I$.
\item The group $G$ acts \emph{proximally} if for every $x, y, z \in (0, 1)$ there exists a sequence $g_i : i \in \N$ such that $x.g_i \to z$ and $y.g_i \to z$.
\end{enumerate}
\end{defi}

A crucial fact, that will be used throughout the paper, is that elements with disjoint support commute.

\begin{rema}
\label{rem:proximal}

Suppose that the diagonal action of $G$ on $(0, 1) \times (0, 1)$ has a dense orbit $O$. Then $G$ acts proximally on $(0, 1)$. Indeed, let $x, y, z \in (0, 1)$, without loss of generality assume $x < y$. Because $O$ is dense, there exists $(x', y') \in O$ such that $x' < x < y < y'$. Using density again, there exists a sequence $g_i$ ($i \in \N$) such that $(x', y').g_i \to (z, z)$. Then $x.g_i \to z$ and $y.g_i \to z$.
\end{rema}

Recall that, for $n\ge 2$, $F_n$ is the group of all piecewise linear orientation preserving homeomorphisms of $[0,1]$ with all slopes powers of $n$ and breakpoints in $\Z[1/n]$. It turns out that $F_n$ admits a standard infinite presentation
\[
F_n=\langle x_0,x_1,\dots\mid x_j^{x_i} = x_{j+(n-1)} \text{ for all } i<j\rangle \text{.}
\]
See for example Proposition~4.8 of~\cite{brown87} (following Propositions~4.1 and 4.4 of~\cite{brown87}, which confirm that we are talking about the same group). Although this presentation (and a related finite presentation) is useful in many contexts, in our case the dynamical definition will be more important.

We introduce the following notation for piecewise linear homeomorphisms:

\begin{defi}
For a piecewise linear element $g \in \Homeo([0, 1])$ and an element $x \in [0, 1]$, we denote by $g'_+(x)$ the right slope at $x$, and by $g'_-(x)$ the left slope at $x$. If the left and right slopes coincide, or if $x \in \{0, 1\}$ and so only one is defined, then we just write $g'(x)$.
\end{defi}


\subsection{Transitivity properties}

The action of $F_n$ on $(0,1)$ stabilizes $(0,1)\cap \Z[1/n]$. However, the action of $F_n$ on $(0,1)\cap \Z[1/n]$ is not transitive when $n > 2$. We refer the reader to~\cite[Section 3]{hydelodha} for a characterization of the orbits; all that we will need is the following:

\begin{cit}[{\cite[Section 3]{hydelodha}}]
\label{cit:proximal}

There exists a dense subset $O \subseteq (0, 1) \cap \Z[1/n]$ that is $F_n$-invariant and on which the action of the commutator subgroup $F_n'$ is highly transitive. That is to say, for any $k\in\N$ and any two tuples $0 < x_1 < \cdots < x_k < 1$ and $0 < y_1 < \cdots < y_k < 1$ with $x_i, y_i \in O$, there exists $g \in F_n'$ such that $x_i.g = y_i$ for all $i$.

In particular, every subgroup $F_n' \leq H \leq F_n$ is proximal on $(0, 1)$.
\end{cit}

The ``in particular'' follows from Remark~\ref{rem:proximal}.


\subsection{Abelianization}

It is clear from the presentation that the abelianization $\Ab(F_n)$ is isomorphic to $\Z^n$. In this section, we will describe this abelianization dynamically. Two characters play a special role:

\begin{defi}
\label{def:chi01}
    Let $\chi_0 \colon F_n \to \Z$ be 	defined as $\chi_0(g) \coloneqq \log_n g'(0)$. Similarly, $\chi_1(g) \coloneqq \log_n g'(1)$. We denote $F_n^c \coloneqq \ker(\chi_0) \cap \ker(\chi_1)$.
\end{defi}

We call $F_n^c$ the \emph{compactly supported subgroup} of $F_n$, since it consists of those homeomorphisms whose support is contained in a compact subset of $(0, 1)$. The quotient $F_n / F_n^c$ is isomorphic to $\Z^2$. This is the end of the story in the important case of Thompson's group $F$, when $n  = 2$. (For $n > 2$, since the abelianization is $\Z^n$, one can find an additional $n-2$ linearly independent characters.)

We will see later that $F_n$ is sharply $\Z^n$-split, i.e., the abelianization map $F_n\to\Z^n$ splits. The following is the key result in this direction. Here an element of $F_n$ is called a \emph{one-bump function} if its support equals a single open interval.

\begin{prop}\label{prop:modelsplitting}
There exist elements $a_0, \ldots, a_{n-1} \in F_n$ such that:
\begin{enumerate}
\item Each $a_i$ is a one-bump function with support $I_i = (\ell_i, r_i)$ for some $\ell_i,r_i$. Moreover, $\ell_0 = 0$ and $\chi_0(a_0)=1$, and $r_{n-1} = 1$ and $\chi_1(a_{n-1})=1$.
\item The $I_i$ are pairwise disjoint, in fact $r_{i-1} < \ell_i$ for $i = 1, \ldots, (n-1)$.
\item The $a_i$ generate the abelianization of $F_n$.
\end{enumerate}
\end{prop}

\begin{proof}
In~\cite{burillo01}, the generators $x_i$ in the standard presentation for $F_n$ are described in terms of the so called tree pair model for elements of $F_n$, and from this it is a straightforward but tedious exercise to view the $x_i$ as piecewise linear homeomorphisms of $[0,1]$. Setting $y_i=x_ix_{i+1}^{-1}$ for each $0\le i\le n-2$, and $y_{n-1}=x_{n-1}^{-1}$, it is easy to see that $y_0,\dots,y_{n-1}$ are one-bump functions, $\chi_0(y_0)=\chi_1(y_{n-1})=1$, and $\{y_0,\dots,y_{n-1}\}$ generates $F_n$. Up to conjugating each $y_i$, and recalling that $F_n$ acts proximally on $(0,1)$ (Citation~\ref{cit:proximal}), we may replace each $y_i$ with an element $a_i$ with support $I_i = (\ell_i, r_i)$ such that $r_0 < \ell_1 < r_1 < \ell_2 < r_2 < \cdots < \ell_{n-2} < r_{n-2} < \ell_{n-1}$. Now the elements $a_0,\dots,a_{n-1}$ satisfy all the desired conditions.
\end{proof}

See Figure~\ref{fig:bumps} for a visualization of the elements $a_i$ constructed in Proposition~\ref{prop:modelsplitting}.

\begin{figure}[htb]
\centering
\begin{tikzpicture}[line width=0.8pt,rotate=45]

\draw[dashed] (0,0) -- (4,0);
\draw (0,0) to[out=40,in=140] (1,0);
\draw (1.5,0) to[out=40,in=140] (2.5,0);
\draw (3,0) to[out=-40,in=-140] (4,0);

\node at (0.5,0.5) {$a_0$};
\node at (2,0.5) {$a_1$};
\node at (3.5,-0.5) {$a_2$};

\end{tikzpicture}
\caption{A graph describing $\{ a_0, a_1, a_2 \}$ in $F_3$, as in Proposition~\ref{prop:modelsplitting}. Note that in order to have $\chi_0(a_0)=\chi_1(a_2)=1$, we need $a_0$ to lie above the $y=x$ line, and $a_2$ to lie below it.}
\label{fig:bumps}
\end{figure}


\subsection{Studying \texorpdfstring{$\Hyp(F_n)$}{Hyp(Fn)} via confining subsets}

In this section, we establish that $F_n$ meets the criteria for the results of Section~\ref{sec:confining} to hold. We start with:

\begin{cit}[{\cite{calegari:PL, commutingconjugates}}]
\label{cit:F_agreeable}
Every subgroup of $F_n$ is agreeable.
\end{cit}

This allows us to exclude non-oriented lineal actions.

\begin{coro}\label{cor:nonoriented:f}
Every lineal action of $F_n$ is oriented.
\end{coro}

\begin{proof}
By Corollary~\ref{cor:nonoriented} and Citation~\ref{cit:F_agreeable}, it suffices to show that $F_n$ does not surject onto $D_\infty$, and this is immediate since every proper quotient of $F_n$ is abelian~\cite[Theorem~4.13]{brown87}.
\end{proof}

We can also exclude general type actions, as a consequence of the following and the well-known Ping-Pong Lemma:

\begin{cit}[\cite{brin85}]
\label{cit:nofree}
    $F_n$ does not contain non-abelian free subgroups. Therefore, $F_n$ does not admit general type actions.
\end{cit}

Finally, we note:

\begin{coro}
\label{cor:Fsplit}
$F_n = F_n' \rtimes \Z^n$ is sharply $\Z^n$-split.
\end{coro}

\begin{proof}
The elements $a_0, \ldots, a_{n-1}$ from Proposition~\ref{prop:modelsplitting} are disjointly supported, and so they commute. This defines a section $\Z^n \to F_n$ of the abelianization map.
\end{proof}

At this point we can write $F_n = F_n' \rtimes \Z^n$, where $\Z^n = \langle a_0, \ldots, a_{n-1} \rangle$ as in Proposition~\ref{prop:modelsplitting}. We fix this $\Z^n$-splitting for the rest of the paper. In particular, note that the restriction of $\chi_0$, respectively $\chi_1$, to this copy of $\Z^n$ maps the product $a_0^{i_0}\cdots a_{n-1}^{i_{n-1}}$ to $i_0$, respectively $i_{n-1}$.


\subsection{Stabilizers in \texorpdfstring{$F_n$}{Fn} and its subgroups}

Later on, we will need some properties of point and interval stabilizers in $F_n$. We first fix the following notation:

\begin{defi}[Fixators, rigid stabilizers]
For an interval $I \subseteq [0, 1]$ (open, closed, or half-open on either side) and a group $G \leq \Homeo([0, 1])$, the \emph{fixator} of $I$ in $G$ is the subgroup
\[
\Fix_G(I) \coloneqq \{ g \in G \mid x.g = x \text{ for all } x \in I \}.
\]
When $I = \{ x \}$ is a degenerate interval, we simply write $\Fix_G(x)$ for $\Fix_G(\{x\})$. The \emph{rigid stabilizer} of $I$ in $G$ is the subgroup $G(I) \coloneqq \left\{ g \in G \mid \overline{\Supp(g)} \subseteq I \right\}$.
\end{defi}

When the interval is written explicitly, we avoid redundancy of parentheses, so for example we write $G[x, y]$ instead of $G([x, y])$. Note that with this notation, we have for example $F_n^c = F_n(0, 1)$.

\begin{defi}[$F_n$-like homeomorphisms]
We say that a homeomorphism between open intervals $I \to J$ is \emph{$F_n$-like} if it is orientation preserving, piecewise linear, with a set of breakpoints that is countable discrete and contained in $\Z[1/n]$, and all slopes powers of $n$.
\end{defi}

Note that when $I = J = (0, 1)$, there are more $F_n$-like homeomorphisms than elements of $F_n$, since in the first case we allow for breakpoints to accumulate at either $0$ or $1$.

\begin{cit}[{\cite[Theorem E16.7]{bieristrebel}}]
\label{cit:bieristrebel}

Let $0 < x < y < 1$ be such that $x, y \in \Z[1/n]$. Then there exists an $F_n$-like homeomorphism $\Phi_{x, y} \colon (x, y) \to (0, 1)$ that conjugates $F_n[x, y]$ to $F_n$.
\end{cit}

Note that the conjugating element really needs to have (countably) infinitely many breakpoints, not merely finitely many; otherwise, it could only intertwine $F_n[x, y]$ and $F_n[0, 1]$ if the orbits were compatible~\cite{hydelodha}.

In the case of rigid stabilizers of open intervals, the isomorphism works even when $x, y$ are not $n$-ary (and moreover, it does not require the previous theorem), as we now show.

\begin{prop}[Rigid stabilizers of open intervals]
\label{prop:rstab:open}

Let $0 < x < y < 1$. Then there exists an $F_n$-like homeomorphism $\Phi_{x, y} \colon (x, y) \to (0, 1)$ that conjugates $F_n(x, y)$ to $F_n^c$.
\end{prop}

\begin{proof}
Recall from Citation~\ref{cit:proximal} that $F_n$ acts transitively on ordered tuples of elements in some dense orbit $O\subseteq (0,1)\cap \Z[1/n]$. Choose a sequence of elements $x < x_i < y_i < y$ such that $x_i, y_i \in O$, the sequence $(x_i)_i$ is strictly decreasing and converges to $x$, and the sequence $(y_i)_i$ is strictly increasing and converges to $y$. Moreover, choose a sequence of elements $0 < p_i < p_{i-1} < q_{i-1} < q_i < 1$ such that $p_i, q_i \in O$ and analogously $p_i \to 0, q_i \to 1$.

Since $F_n$ is transitive on ordered pairs in $O$, $F_n$ can take each $x_i<y_i$ to $p_i<q_i$. In this way, we get $F_n$-like homeomorphisms $\Phi_i \colon (x_i, y_i) \to (p_i, q_i)$ that conjugate $F_n(x_i, y_i)$ to $F_n(p_i, q_i)$ simply by conjugating by an appropriate element of $F_n$. Moreover, by defining the elements of $F_n$ inductively and by pieces, we may ensure that each $\Phi_i$ extends $\Phi_{i-1}$.

This defines an orientation preserving homeomorphism $\Phi \colon (x, y) \to (0, 1)$ that conjugates the directed union of the groups $F_n(x_i, y_i)$, namely $F_n(x, y)$, to the directed union of the $F_n(p_i, q_i)$, namely $F_n(0, 1) = F_n^c$. In fact, each of the $\Phi_i$ is constructed via conjugation in $F_n$, and so has finitely many breakpoints. It follows that the set of breakpoints of $\Phi$ is countable and discrete in $(x, y)$, and so $\Phi$ is $F_n$-like.
\end{proof}

Proposition~\ref{prop:rstab:open} shows that the behavior of rigid stabilizers of open intervals is independent of the nature of the endpoints. However, the behavior of stabilizers of points does depend on the nature of the points, with different behavior when a fixed point is $n$-ary, rational but not $n$-ary, or irrational. We will now describe the behavior in each of the three cases; when $n = 2$ this was already worked out in~\cite{golansapir}.

The $n$-ary case is immediate, because of the piecewise nature of $F_n$:

\begin{lemm}[$n$-ary fixed points]
\label{lem:fix:nary}

Let $t \in \Z[1/n] \cap (0, 1)$, Then $\Fix_{F_n}(t)$ decomposes as $F_n[0, t] \times F_n[t, 1]$.\qed
\end{lemm}

Note however that this product decomposition is not always stable under replacing $F_n$ in the notation with a subgroup. For example using $F_n'$, under this decomposition $\Fix_{F_n'}(t)$ corresponds to those elements $(g, h)$ of $F_n[0, t] \times F_n[t, 1]$ such that $\Ab(g) = - \Ab(h)$, and it is easy to see that $\Fix_{F_n'}(t)$ is identified with a subgroup of $F_n(0, t] \times F_n[t, 1)$ whose projection onto each factor is surjective. In particular, $F_n'(t)$ does not decompose as $F_n'[0,t]\times F_n'[t,1]$.

The irrational case is also relatively straightforward:

\begin{lemm}[Irrational fixed points]
\label{lem:fix:irrational}

Let $t \in (0, 1) \setminus \Q$. Then every element $g \in F_n$ that fixes $t$ also fixes a neighborhood of $t$. Thus $\Fix_{F_n}(t)$ decomposes as $F_n[0, t] \times F_n[t, 1] = F_n[0, t) \times F_n(t, 1]$.
\end{lemm}

\begin{proof}
Suppose that, at $t$, the element $g$ takes the form $x \mapsto n^k x + y$, for $k \in \Z$ and $y \in \Z[1/n]$. If $k \neq 0$, then $t = t.g$ gives $t = y/(1-n^k) \in \Q$. Since $t$ is irrational, we conclude that $k=0$. Moreover, $t$ cannot be a break point, so the slope of $g$ has to be identically $1$ in a neighborhood of $t$, i.e., $g$ must fix a neighborhood of $t$. The decomposition into a direct product is now immediate.
\end{proof}

Finally there is the case when the fixed point is rational but not $n$-ary, which is more delicate:

\begin{lemm}[Non $n$-ary rational fixed points]
\label{lem:fix:rational}

Let $t \in (0, 1) \cap \Q \setminus \Z[1/n]$. Then the map $\Fix_{F_n}(t) \to \Z$ given by $g \mapsto \log_n g'(t)$ is a well defined non-trivial homomorphism. The image only depends on the $F_n$-orbit of $t$. Thus $\Fix_{F_n}(t)$ fits into a short exact sequence
\[1 \to F_n[0, t) \times F_n(t, 1] \to \Fix_{F_n}(t) \to \Z \to 1;\]
where the element of $\Fix_{F_n}(t)$ mapping to the generator of the image in $\Z$ may be chosen to belong to $F_n'$ and to be supported on an arbitrarily small open neighborhood of $t$.
\end{lemm}

\begin{proof}
Since $t \notin \Z[1/n]$, it cannot be a breakpoint, so $g'(t)$ is well defined for all $g \in F_n$, and it is easily seen that we get the homomorphism in the statement, by applications of the chain rule and log rules.
To see that this homomorphism is non-trivial, write $t = p/q$, where $p$ and $q$ are coprime, and so by hypothesis $q$ is not a power of $n$. We multiply the numerator and denominator suitably to obtain an expression $t = p'/q' n^\ell$, where $q'$ is coprime to $n$, and greater than $1$. Let $o\ge 1$ be the order of $n$ in $(\Z/q'\Z)^\times$; so $q'$ divides $(n^o - 1)$. By multiplying the numerator and denominator again, we obtain an expression
\[
t = \frac{m}{(n^o - 1) n^\ell} = \frac{1}{n^o-1} \frac{m}{n^\ell}\text{.}
\]
Setting $y = m/n^\ell \in \Z[1/n]$, we obtain the expression $t = n^o t - y$.

Let $I$ be an open interval containing $t$ such that the map $x \mapsto n^o x - y$, which fixes~$t$, maps $I$ inside a compact subset of $(0, 1)$. This can then be completed to an element of $F_n$, and even to an element of $F_n'$, up to multiplying by suitable conjugates of the elements $a_i$ from Proposition~\ref{prop:modelsplitting} supported away from $t$. Since $o \ge 1$, this shows that the homomorphism in the statement is non-trivial, with image generated by an element of $F_n'$, and the kernel is easily identified.

Finally, let $g \in \Fix_{F_n}(t)$ and let $c \in F_n$. Then $g^c \in \Fix_{F_n}(t.c)$, and moreover $t.c \in \Q \setminus \Z[1/n]$ and the slope of $g$ at $t$ coincides with the slope of $g^c$ at $t.c$. This shows that the image of $\Fix_{F_n}(t) \to \Z$ given by $g \mapsto \log_n g'(t)$ does not change if we replace $t$ by a rational in the same orbit, and concludes the proof.
\end{proof}

\begin{rema}
The proof actually shows that the generator of the image of $\log_n g'(t)$ coincides with the order of $n$ modulo $q'$, where $q'$ is the largest factor of the denominator of $t$ that is coprime to $n$. The invariance over $F_n$-orbits is also reflected in this property~\cite{hydelodha}.
\end{rema}

\section{Global structure of \texorpdfstring{$\Hyp(F_n)$}{Hyp(Fn)}}\label{sec:global}

In this section we analyze confining subsets in $F_n$, and find that for ``almost all'' $\rho$ there are no quasi-parabolic actions of $F_n$ with Busemann pseudocharacter equivalent to $\rho$. More precisely, this happens for any $\rho\ne \chi_0,\chi_1$. Since $F_n$ is agreeable (Citation~\ref{cit:F_agreeable}) and sharply $\Z^n$-split (Corollary~\ref{cor:Fsplit}), we can use the results from Section~\ref{sec:confining}, and so really what we are proving is that for any $\rho\ne \chi_0,\chi_1$, the set $\Conf^\ast_\rho(F_n)$ is empty. Here $\Conf^\ast_\rho(F_n)$ is as in Definition~\ref{def:conf_poset}. Throughout this, we fix the sharp $\Z^n$-splitting $F_n = F_n' \rtimes \mathbb{Z}^n$ as in Proposition~\ref{prop:modelsplitting}, where $\Z^n$ is the subgroup generated by $\{a_0,\dots,a_{n-1}\}$. In particular, each $a_i$ is a one-bump function with support $I_i = (\ell_i, r_i)$, these intervals are pairwise disjoint, and we have $\ell_0 = 0, r_{n-1} = 1$.

We should briefly remark that there is a connection between hyperbolic actions and Bieri--Neumann--Strebel (BNS) invariants of groups, as in~\cite{bieri87}; see for example~\cite[Subsection~4.3]{confining}. The BNS-invariants of the $F_n$ have been completely determined~\cite[Theorem~8.1]{bieri87}, and thanks to an equivalent characterization of the BNS-invariants due to Brown~\cite[Corollary~7.4]{brown87trees}, it is clear that there are no $\rho$-confining subsets for $\rho \neq \chi_0, \chi_1$ whose corresponding hyperbolic actions are on $\R$-trees. This does not immediately help us here though, because $\R$-trees are just a special case of hyperbolic spaces, so the connection to BNS-invariants is more just ``in spirit''.

\begin{prop}\label{prop:usually_empty_conf}
Let $\rho$ be a character of $\Z^n$ such that there exist $i\ne j$ with $\rho(a_i)\ne 0$ and $\rho(a_j)\ne 0$. Then $\Conf^\ast_\rho(F_n)=\emptyset$.
\end{prop}

\begin{proof}
For each $k$, note that the centralizer $C_{F_n'}(a_k)$ contains the subgroup $\Fix_{F_n'}(I_k)$ of all elements of $F_n'$ fixing $I_k$. By Proposition~\ref{prop:empty_conf}, it therefore suffices to prove that $F_n'$ is boundedly generated by $\Fix_{F_n'}(I_i)\cup \Fix_{F_n'}(I_j)$. We claim that every element of $F_n'$ is a product of at most three elements of this union.
Let $f\in F_n'$, so in particular $f \in F_n^c$ and so the support of $f$ is contained in $(\varepsilon,\delta)$ for some $0<\varepsilon<\delta<1$. Since $i\ne j$ we know that $I_i$ and $I_j$ are disjoint. Say without loss of generality that $i<j$. If $\varepsilon>r_i$ then $f\in \Fix_{F_n'}(I_i)$ and we are done, so suppose $\varepsilon\le r_i$. Since $I_i$ and $I_j$ are disjoint, we can choose some element $f'\in \Fix_{F_n'}(I_j)$ such that $\varepsilon.f'>r_i$. Now the support of $f^{f'}$ is contained in $(r_i,1)$ and hence $f^{f'}\in\Fix_{F_n'}(I_i)$. We conclude that
\[
f\in f'\Fix_{F_n'}(I_i)(f')^{-1} \subseteq \Fix_{F_n'}(I_j)\Fix_{F_n'}(I_i)\Fix_{F_n'}(I_j)\text{,}
\] and we are done.
\end{proof}

At this point we are free to assume that $\rho(a_i)\ne 0$ for just a single $i$. Our next goal is to prove that in fact $\Conf^\ast_\rho(F_n)=\emptyset$ also holds if this $i$ is anything other than $0$ or $n-1$.

\begin{prop}\label{prop:middle_empty_conf}
Let $\rho$ be a character of $\Z^n$ such that $\rho(a_i)\ne 0$ for some $0<i<n-1$, and $\rho(a_j)=0$ for all $j\ne i$. Then $\Conf^\ast_\rho(F_n)=\emptyset$.
\end{prop}

\begin{proof}
For notational ease, let us write $a$ for $a_i$, $I$ for the interval $I_i$, and $\ell$ and $r$ for the endpoints $\ell_i$ and $r_i$ of $I$. Up to replacing $a$ by its inverse, let us assume that $\rho(a) > 0$. Let $Q\subseteq F_n'$ be a $\rho$-confining subset. We need to prove that $Q=F_n'$. Since $0<i<n-1$, we know that $(r_{i-1},\ell)\ne\emptyset$ and $(r,\ell_{i+1})\ne\emptyset$. By Proposition~\ref{prop:empty_conf}, we know that $C_{F_n'}(a) \subseteq Q$, and so $\Fix_{F_n'}(I)\subseteq Q$. 

We claim that there exists $q_\ell\in Q$ such that $x_\ell.q_\ell<\ell$ for some $x_\ell\in (\ell,r)$. First let $q_\ell'\in F_n'$ be any element supported on $(r_{i-1},r)$, such that $\ell.q_\ell'<\ell$, so then we also get $x_\ell'.q_\ell'<\ell$ for some $x_\ell'\in (\ell,r)$. By Lemma~\ref{lem:conf:cyclic}, we can choose $k \geq 0$ such that $q_\ell\coloneqq (q_\ell')^{a^k}\in Q$. Let $x_\ell\coloneqq x_\ell'.a^k$. Then $x_\ell\in (\ell,r)$ and
\[
x_\ell.q_\ell = (x_\ell'.a^k).(q_\ell')^{a^k} = (x_\ell'.q_\ell').a^k = x_\ell'.q_\ell'<\ell
\]
since $a^k$ fixes $[0,\ell]$. A symmetric argument produces an element $q_r\in Q$ and $x_r\in(\ell,r)$ such that $x_r.q_r>r$.

By now we know that $Q$ contains $\Fix_{F_n'}(I)\subseteq Q$ and contains the elements $q_\ell$ and $q_r$ constructed above. Now we are ready to begin proving that $Q=F_n'$, for which it suffices to show that $F_n'$ is bounded in the generating set $Q$. Let us assume that $x.a>x$ for all $x$ in the support $I$ of $a$; since $a$ is a one-bump function, the alternative is that $x.a<x$ for all $x\in I$, and this case works analogously. Let $f\in F_n'$, so in particular $f \in F_n^c$ and so $f$ is supported on $(\varepsilon,\delta)$ for some $0<\varepsilon<\delta<1$. If $\delta\le \ell$ then $f\in \Fix_{F_n'}(I)$, so $f\in Q$ and we are done. Now suppose $\ell < \delta < r$. Since $a$ is a one-bump function satisfying $x.a>x$ for all $x$ in its support $(\ell,r)$, we can choose $k\ge 0$ such that $\delta.a^{-k}<x_\ell$, and so $f^{a^{-k}q_\ell}\in \Fix_{F_n'}(I) \subseteq Q$. This implies that $f \in Q^{q_\ell^{-1} a^k} \subseteq (Q^3)^{a^k} \subseteq Q^3$, where the last inclusion is given by \StayInConf. Finally, suppose $r \le \delta$. Choose $f'\in \Fix_{F_n'}(I)\subseteq Q$ such that $\delta.f'<x_r.q_r$, so $\delta.f'q_r^{-1}<r$. Like in the previous case, we can now choose $k\ge 0$ such that $\delta.f'q_r^{-1} a^{-k}<x_\ell$, so $f^{f'q_r^{-1} a^{-k}q_\ell}\in  \Fix_{F_n'}(I) \subseteq Q$, and so $f\in Q^{q_\ell^{-1} a^k q_r (f')^{-1}} \subseteq Q^7$. We conclude that in all cases $f\in Q^7$, which concludes the proof.
\end{proof}

This leaves us with only four characters of $\Z^n$ to consider up to equivalence, since precisely one of $a_0$ or $a_{n-1}$ can have a positive or negative character value, and the other must be sent to zero. In fact by Proposition~\ref{prop:modelsplitting} these four characters of $\Z^n$ are equivalent to the four characters $\pm\chi_0$ and $\pm\chi_1$ of $F_n$, restricted to $\Z^n$. Our next result is that $\Conf^\ast_{-\chi_0}(F_n)$ and $\Conf^\ast_{-\chi_1}(F_n)$ are also empty, leaving $\chi_0$ and $\chi_1$ as the only candidates for yielding quasi-parabolic actions. (Then later we will show that these actually yield an incredible number of quasi-parabolic actions.)

\begin{prop}\label{prop:last_empty_conf}
The sets $\Conf^\ast_{-\chi_0}(F_n)$ and $\Conf^\ast_{-\chi_1}(F_n)$ are empty.
\end{prop}

\begin{proof}
We will prove $\Conf^\ast_{-\chi_0}(F_n)=\emptyset$, and the other statement is analogous. Let $Q\subseteq F_n'$ be $(-\chi_0)$-confining. By Proposition~\ref{prop:empty_conf}, we know that $C_{F_n'}(a_0) \subseteq Q$, and so $\Fix_{F_n'}(I_0)=F_n'[r_0,1)\subseteq Q$. By the same argument as in the proof of Proposition~\ref{prop:middle_empty_conf}, $Q$ also contains an element $q$ such that $x.q>r_0$ for some $x\in (0,r_0)$. Now we are ready to show that $Q=F_n'$, for which it suffices to show that $F_n' \subseteq Q^3$. Let $f\in F_n'$, say supported on $(\varepsilon,\delta)$ for some $0<\varepsilon<\delta<1$. Choose $k\ge 0$ such that $\varepsilon.a_0^k>x$, so $\varepsilon.a_0^k q > r_0$. Now $f^{a_0^k q}\in F_n'[r_0,1)\subseteq Q$, so $f\in (Q^3)^{a_0^{-k}}$. Since $(-\chi_0)(a_i^{-k})=k\ge 0$, we conclude that $f\in Q^3$ by \StayInConf, which concludes the proof.
\end{proof}

We now have everything in place to give a global description of $\Hyp(F_n)$. We refer the reader to Figure~\ref{intro:fig:global} in the introduction for a picture.

\begin{theo}[Theorem~\ref{intro:thm:global}]\label{thm:global}
The poset $\Hyp(F_n)$ has the following structure.
\begin{enumerate}
   \item $\Hyp_{gt}(F_n) = \emptyset$.
   \item $\Hyp_{qp}(F_n)$ consists of two incomparable isomorphic join semilattices $\X_0$ and $\X_1$.
   \item $\Hyp_{\ell}(F_n) = \Hyp_{\ell}^+(F_n)$ is an uncountable antichain, with two distinguished elements $[\pm\chi_0]$ and $[\pm\chi_1]$ sitting under $\X_0$ and $\X_1$, respectively.
   \item $\Hyp_e(F_n)$ consists of the unique smallest element.
\end{enumerate}
\end{theo}

\begin{proof}
The last item is always true. The fact that $\Hyp_{gt}(F_n) = \emptyset$ is the content of Citation~\ref{cit:nofree}. The fact that $\Hyp_{\ell}(F_n) = \Hyp_{\ell}^+(F_n)$ is the content of Corollary~\ref{cor:nonoriented:f}. This is an uncountable antichain by Proposition~\ref{prop:lineal}, with two distinguished elements corresponding to $[\pm\chi_0]$ and $[\pm\chi_1]$.

We define $\X_0 \coloneqq \Hyp_{qp}^{\chi_0}(F_n)$, and similarly $\X_1 \coloneqq \Hyp_{qp}^{\chi_1}(F_n)$. They are join semilattices by Proposition~\ref{prop:join}. They are isomorphic because the automorphism $\Homeo[0, 1] \to \Homeo[0, 1]$ given by conjugation by $x \mapsto (1 - x)$ preserves $F_n$ and conjugates $\chi_0$ to~$\chi_1$. By construction, they sit above $[\pm\chi_0]$ and $[\pm\chi_1]$, respectively, and above no other lineal actions by Proposition~\ref{prop:linvsqp}. Finally, that $\Hyp_{qp}(F_n)$ consists only of $\X_0$ and $\X_1$ follows from the propositions proved above in this section that show no other character admits strictly confining subsets.
\end{proof}
